% Insertable Q1 chapter. Compile from paper/; preview shell: q1-preview.tex.
% Numerical tables and macros are generated by src/q1_chapter_assets.py.
% Values derived from accepted Q1 authority and bound audit data.
\newcommand{\QOneTCentre}{33.5753}
\newcommand{\QOneTSurface}{36.7856}
\newcommand{\QOneTGap}{3.2102}
\newcommand{\QOneTEnvironment}{41.513}
\newcommand{\QOneCCentre}{2.5500}
\newcommand{\QOneCSurface}{1.5102}
\newcommand{\QOneCCentreFine}{2.549992404}
\newcommand{\QOneCeEnd}{0.03307}
\newcommand{\QOneAlpha}{1.68856\times10^{-7}}
\newcommand{\QOneDInitial}{4.93766\times10^{-9}}
\newcommand{\QOneBiT}{1.3889}
\newcommand{\QOneHeatTime}{2368.9}
\newcommand{\QOneMassTime}{81010.1}
\newcommand{\QOneHeatLength}{17.43}
\newcommand{\QOneMassLength}{2.98}
\newcommand{\QOneDLastSurface}{3.88296\times10^{-9}}
\newcommand{\QOneTEstimate}{2.38308\times10^{-8}}
\newcommand{\QOneCEstimate}{1.50463\times10^{-7}}
\newcommand{\QOneTIndependent}{3.06892\times10^{-7}}
\newcommand{\QOneCIndependent}{2.44260\times10^{-7}}

\section{问题1：预热平衡阶段的模型建立与求解}
\label{sec:q1}

\subsection{问题分析与建模范围}

问题1要求根据烘房环境随时间的变化，计算前30分钟内药材各径向位置的温度与水分浓度。药材近似为长$L=0.25\,\mathrm{m}$、半径$R=0.02\,\mathrm{m}$的圆柱。由于所求结果具有空间分辨要求，且表面交换存在有限阻力，采用内部均匀的集中参数模型会丢失中心与表面的差异。本问建立\textbf{半解析传热与非线性守恒传质的径向组合模型}：利用常物性导热方程的线性结构求解温度，以守恒通量离散保留含水率对扩散系数的非线性影响。

模型针对圆柱中部代表截面，假设材料在周向均匀，忽略该截面的轴向传输；在预热阶段保持半径和长度不变。以题定物性估算，1800 s内的热扩散尺度约为$\QOneHeatLength\,\mathrm{mm}$，初始水分扩散尺度约为$\QOneMassLength\,\mathrm{mm}$，均小于中截面至端面的125 mm距离。这为中部径向近似提供尺度依据，但不将其扩展为端部局部场的预测。采用附录2的常物性，不引入蒸发潜热项；水分扩散系数仅依赖局部含水率。因此，当前两场共用几何及环境输入，可分别求解，并未构成温度与水分的双向强耦合。

\subsection{控制方程与初边值条件}

设$t$为时间（s），$r$为到中心轴的距离（m）；$T(r,t)$为药材温度（℃），$C(r,t)$为题定干基含水率（kg/kg）。$T_a(t)$为烘房温度，$C_e(t)$为表面传质条件中的等效环境驱动，$T_s(t)=T(R,t)$、$C_s(t)=C(R,t)$分别为药材表面状态。计算内部统一采用m；题目要求的结果表以cm列示距离。温差的℃数值与K一致，不改变输入温度的数值基准。

附录2给出的密度$\rho$、比热容$c_p$、导热系数$k$、表面对流换热系数$h_T$与传质系数$h_C$为
\begin{equation}
\begin{aligned}
\rho&=820\,\mathrm{kg/m^3},&
c_p&=2600\,\mathrm{J/(kg\cdot K)},&
k&=0.36\,\mathrm{W/(m\cdot K)},\\
h_T&=25\,\mathrm{W/(m^2\cdot K)},&
h_C&=8\times10^{-7}\,\mathrm{m/s}.&&
\end{aligned}
\label{eq:q1-parameters}
\end{equation}
在圆柱壳层内分别建立显热与题定含水率的积分平衡，令壳层厚度趋于零，得到
\begin{equation}
\begin{aligned}
\rho c_p\frac{\partial T}{\partial t}
&=\frac{k}{r}\frac{\partial}{\partial r}
 \left(r\frac{\partial T}{\partial r}\right),\\
\frac{\partial C}{\partial t}
&=\frac{1}{r}\frac{\partial}{\partial r}
 \left[rD(C)\frac{\partial C}{\partial r}\right],
\qquad 0<r<R,\quad 0<t\le1800.
\end{aligned}
\label{eq:q1-pde}
\end{equation}
其中题定经验扩散系数为
\begin{equation}
D(C)=D_0\exp\!\left(-\frac{a}{C}\right),
\qquad D_0=7\times10^{-9}\,\mathrm{m^2/s},\quad a=0.89\,\mathrm{kg/kg}.
\label{eq:q1-diffusion}
\end{equation}
经验式中的$C$直接使用kg/kg口径的题定数值，不换成百分数、湿基含水率或体积浓度。各处扩散系数随当地$C$计算，不以环境量或表面量代替内部状态。

初始状态在空间上均匀；中心满足轴对称条件，表面保留内外状态差所驱动的有限交换：
\begin{equation}
\begin{aligned}
T(r,0)&=28\,{}^{\circ}\mathrm C,&
C(r,0)&=2.55\,\mathrm{kg/kg},\qquad 0\le r\le R,\\
T_r(0,t)&=0,& C_r(0,t)&=0,\qquad t>0,\\
-kT_r(R,t)&=h_T[T_s(t)-T_a(t)],&
-D(C_s)C_r(R,t)&=h_C[C_s(t)-C_e(t)],\qquad t>0.
\end{aligned}
\label{eq:q1-ibc}
\end{equation}
上述符号以径向向外为正：加热时$T_s<T_a$，热量向内传递；干燥时$C_s>C_e$，水分向外传递。$h_C$保留题定m/s单位，不额外乘除密度。

附件1记录的是烘房环境，而非药材内部观测。本问取其温度列为$T_a$，并\textbf{假设水分列可作为等效$C_e$}用于式\eqref{eq:q1-ibc}。这一映射提供本题的边界闭合，不表示已经获得材料与空气间的平衡含水率关系。对0—1800 s内的31个真实采样节点$t_j=60j\,\mathrm{s}$，环境量$X\in\{T_a,C_e\}$按
\begin{equation}
X(t)=X(t_j)+\frac{t-t_j}{t_{j+1}-t_j}
       [X(t_{j+1})-X(t_j)],\qquad t_j\le t\le t_{j+1}
\label{eq:q1-environment}
\end{equation}
逐段线性给定，保持全部节点原值。附件2的半径变化数据不用于本问固定半径模型，附件3仅提供结果工作簿模板。

初始烘房温度为28 ℃，与温度初值相容；初始等效$C_e=0.01963\,\mathrm{kg/kg}$，与均匀药材初值对应的零梯度不满足同一时刻的传质交换条件。计算保留题定均匀初态，在$t>0$施加Robin边界，由扩散形成启动边界层，不通过修改初始表面值消除不相容性。

\subsection{半解析温度解与非线性水分解}

\subsubsection{圆柱模态展开与环境分段推进}

令热扩散率$\alpha=k/(\rho c_p)$（$\mathrm{m^2/s}$）、无量纲半径$x=r/R$及相对环境温差$\theta=T-T_a(t)$（℃）。变换后表面为齐次Robin条件，内部驱动为$-T_a'(t)$。圆柱径向特征方程的中心正则解为零阶第一类Bessel函数$J_0$\cite{q1-dlmf}，其正特征根按升序满足
\begin{equation}
\lambda_nJ_1(\lambda_n)=Bi_TJ_0(\lambda_n),
\qquad Bi_T=\frac{h_TR}{k}=\QOneBiT.
\label{eq:q1-roots}
\end{equation}
这里$J_1$为一阶第一类Bessel函数。$Bi_T$并不小，故内部热阻及表面温差均予保留。以圆柱权重$x$投影常函数1，定义
\begin{equation}
a_n=\frac{\int_0^1xJ_0(\lambda_nx)\,\mathrm dx}
           {\int_0^1xJ_0^2(\lambda_nx)\,\mathrm dx}
   =\frac{2J_1(\lambda_n)}
           {\lambda_n[J_0^2(\lambda_n)+J_1^2(\lambda_n)]},
\qquad \beta_n=\frac{\alpha\lambda_n^2}{R^2}.
\label{eq:q1-projection}
\end{equation}
其中$a_n$无量纲，$\beta_n$为模态衰减率（$\mathrm{s^{-1}}$）。温度用$M=1024$个模态近似表示为
\begin{equation}
T_M(r,t)=T_a(t)+\sum_{n=1}^{M}b_n(t)J_0(\lambda_nr/R),
\qquad b_n'+\beta_nb_n=-a_nT_a'(t).
\label{eq:q1-temperature}
\end{equation}
初始投影$b_n(0)=a_n[28-T_a(0)]=0$。在第$j$个环境线性段，记斜率$s_j=[T_a(t_{j+1})-T_a(t_j)]/60$，则对$0\le\tau\le60\,\mathrm{s}$有
\begin{equation}
b_n(t_j+\tau)=e^{-\beta_n\tau}b_n(t_j)
 -a_ns_j\frac{1-e^{-\beta_n\tau}}{\beta_n}.
\label{eq:q1-mode-step}
\end{equation}
式\eqref{eq:q1-mode-step}在每段内解析推进模态，并以段末状态衔接下一段，保证温度连续；小指数差采用稳定的指数差函数求值。该处理消除了有限模态系统的时间离散误差，剩余的模态截断、特征根与浮点求值误差另行检验。

\subsubsection{守恒环带离散、表面状态与时间积分}

为保留$D(C)$的局部非线性，引入通量势
\begin{equation}
\Phi(C)=\int_0^CD(s)\,\mathrm ds
       =D_0C E_2(a/C),\qquad
E_2(z)=\int_1^\infty\frac{e^{-zu}}{u^2}\,\mathrm du,\quad C>0.
\label{eq:q1-potential}
\end{equation}
由$\Phi'(C)=D(C)>0$，扩散项可写为$\frac1r(r\Phi_r)_r$，且势差与含水率差同号，保留从高含水率向低含水率的通量方向。$\Phi$并非$D(C)C$；当两状态接近时，直接对区间上的$D$进行12点Gauss积分求势差，以减小相近势值相减的消减误差。

将半径等分为$N=32768$个环带，$\Delta r=R/N$。第$i$环（$i=0,\ldots,N-1$）的边界为$r_{i-1/2}=i\Delta r$、$r_{i+1/2}=(i+1)\Delta r$，代表半径为$\hat r_i=(i+1/2)\Delta r$。统一取$L=0.25\,\mathrm m$，环体积、面面积与未知量分别为
\begin{equation}
\begin{aligned}
V_i&=\pi L(r_{i+1/2}^2-r_{i-1/2}^2),\qquad A_f=2\pi Lr_f,\\
\bar C_i(t)&=\frac{2\pi L}{V_i}
 \int_{r_{i-1/2}}^{r_{i+1/2}}rC(r,t)\,\mathrm dr.
\end{aligned}
\label{eq:q1-cell}
\end{equation}
$\bar C_i$为环带体积平均值，并非代表半径处的精确点值。对内部面构造共用的向外通量
\begin{equation}
\begin{aligned}
F_{i+1/2}&=\frac{A_{i+1/2}}{\hat r_{i+1}-\hat r_i}
 [\Phi(\bar C_i)-\Phi(\bar C_{i+1})],\qquad i=0,\ldots,N-2,\\
V_i\frac{\mathrm d\bar C_i}{\mathrm dt}&=F_{i-1/2}-F_{i+1/2}.
\end{aligned}
\label{eq:q1-fv}
\end{equation}
相邻环使用同一面通量并取相反符号，使内部交换在全域求和时抵消。中心面的面积为零，令$F_{-1/2}=0$，无需在$r=0$计算$1/r$。以环均值代入势差属于空间近似，其一致性和点值精度由后述空间细化验证。

外表面另设独立代数状态$C_s$。取末环代表位置至表面的距离$\delta=R-\hat r_{N-1}=\Delta r/2$，连接内部扩散与表面交换：
\begin{equation}
g(C_s)=\frac{\Phi(C_s)-\Phi(\bar C_{N-1})}{\delta}
          +h_C(C_s-C_e)=0,\qquad
F_{N-1/2}=A_Rh_C(C_s-C_e).
\label{eq:q1-surface}
\end{equation}
其中$A_R=2\pi LR$。由于$g'(C_s)=D(C_s)/\delta+h_C>0$，正值条件下根位于$C_e$与$\bar C_{N-1}$之间且唯一，可用括区间的Brent法求解，绝对根容差取$5\times10^{-15}\,\mathrm{kg/kg}$并检查原始边界残差。根的唯一性不消除半环连接的空间误差，也不允许以末环均值替代$C_s$。

空间半离散系统采用三阶段、五阶\textbf{Radau IIA隐式积分}，相对容差$10^{-12}$、绝对容差$10^{-14}$，最大步长0.5 s。其三阶嵌入估计以加权均方根范数控制局部误差，不能直接解释为全时段每点误差上界；整数秒状态使用三次配点稠密输出恢复\cite{q1-radau}。积分准确到达每个60 s环境节点，以完整环状态继续下一段；初步长按$0.01\Delta r^2/D(2.55)$设置，以分辨启动边界层的快速变化。输出间隔1 s不等于实际积分步长。

由环均值恢复题定点值时，中心利用偶对称二次局部展开得到
\begin{equation}
C(0,t)\approx\frac{5\bar C_0(t)-\bar C_1(t)}4.
\label{eq:q1-center}
\end{equation}
其他内部半径以$r^2$为变量，用三个邻近环的精确积分矩匹配二次多项式后求点值；表面直接输出式\eqref{eq:q1-surface}求得的$C_s$。由此区分环均值、内部点值和几何表面的模型状态，重构误差与空间离散一并在相同物理坐标处比较。开发阶段的后向Euler仅保留为基线，未用于最终结果表；更细的65536环也仅用于参照核验。

\subsection{计算结果与机理解释}

在$t=100,300,600,900,1200,1500,1800\,\mathrm s$及$r=0,0.5,1,1.5,2\,\mathrm{cm}$提取结果，分别得到表\ref{tab:q1-temperature}和表\ref{tab:q1-moisture}。两表均由同一份全精度解统一舍入生成；每秒、每0.1 cm的完整结果保存于\texttt{result1.xlsx}，两工作表各包含1800$\times$21个数值，时间从1 s开始。

% Generated from frozen q1-authority.npz; do not hand-edit numeric cells.
% SHA-256: 8aab6a875752f60bb798519d60e2532a54585874b92e18183adb861571d320a4
\begin{table}[!htbp]
\centering
\small
\caption{30分钟内药材的温度（单位：℃）}
\label{tab:q1-temperature}
\setlength{\tabcolsep}{9pt}
\renewcommand{\arraystretch}{1.12}
\begin{tabular*}{0.86\textwidth}{@{\extracolsep{\fill}}rccccc}
\toprule
时间/s & \multicolumn{5}{c}{到药材中心的距离/cm} \\
\cmidrule(lr){2-6}
 & 0 & 0.5 & 1 & 1.5 & 2 \\
\midrule
100 & 28.0001 & 28.0003 & 28.0040 & 28.0327 & 28.1801 \\
300 & 28.0408 & 28.0635 & 28.1514 & 28.3681 & 28.8489 \\
600 & 28.4533 & 28.5360 & 28.8039 & 29.3159 & 30.1652 \\
900 & 29.3243 & 29.4583 & 29.8755 & 30.6161 & 31.7304 \\
1200 & 30.5427 & 30.7098 & 31.2223 & 32.1126 & 33.4276 \\
1500 & 31.9957 & 32.1867 & 32.7660 & 33.7463 & 35.1203 \\
1800 & 33.5753 & 33.7720 & 34.3642 & 35.3621 & 36.7856 \\
\bottomrule
\end{tabular*}
\end{table}

% Generated from frozen q1-authority.npz; do not hand-edit numeric cells.
% SHA-256: 8aab6a875752f60bb798519d60e2532a54585874b92e18183adb861571d320a4
\begin{table}[!htbp]
\centering
\small
\caption{30分钟内药材的水分浓度（干基，单位：kg/kg）}
\label{tab:q1-moisture}
\setlength{\tabcolsep}{9pt}
\renewcommand{\arraystretch}{1.12}
\begin{tabular*}{0.86\textwidth}{@{\extracolsep{\fill}}rccccc}
\toprule
时间/s & \multicolumn{5}{c}{到药材中心的距离/cm} \\
\cmidrule(lr){2-6}
 & 0 & 0.5 & 1 & 1.5 & 2 \\
\midrule
100 & 2.5500 & 2.5500 & 2.5500 & 2.5500 & 2.2470 \\
300 & 2.5500 & 2.5500 & 2.5500 & 2.5492 & 2.0508 \\
600 & 2.5500 & 2.5500 & 2.5500 & 2.5353 & 1.8770 \\
900 & 2.5500 & 2.5500 & 2.5497 & 2.5045 & 1.7547 \\
1200 & 2.5500 & 2.5500 & 2.5482 & 2.4646 & 1.6586 \\
1500 & 2.5500 & 2.5499 & 2.5445 & 2.4206 & 1.5787 \\
1800 & 2.5500 & 2.5497 & 2.5383 & 2.3755 & 1.5102 \\
\bottomrule
\end{tabular*}
\end{table}


表\ref{tab:q1-temperature}显示，所列时刻表面升温快于中心。1800 s时中心与表面分别为$\QOneTCentre\,{}^{\circ}\mathrm C$和$\QOneTSurface\,{}^{\circ}\mathrm C$，按全精度解计算的径向温差约为$\QOneTGap\,{}^{\circ}\mathrm C$；同期烘房为$\QOneTEnvironment\,{}^{\circ}\mathrm C$，表面尚未达到环境温度。由题定参数得到$\alpha=\QOneAlpha\,\mathrm{m^2/s}$，径向热扩散时间$R^2/\alpha\approx\QOneHeatTime\,\mathrm s$，与观察时段同量级。因此，在持续变化的环境和有限换热阻力下，热量向内传播需要时间，“预热平衡阶段”不能直接解释为全截面已经达到稳态。

\begin{samepage}
表\ref{tab:q1-moisture}表明，失水主要集中于外层：1800 s时表面为$\QOneCSurface\,\mathrm{kg/kg}$，中心仍显示为$\QOneCCentre\,\mathrm{kg/kg}$。中心的全精度计算值约为$\QOneCCentreFine\,\mathrm{kg/kg}$，应理解为变化很小、四位小数下与初值相同，而非严格不变。初始扩散系数$D(2.55)=\QOneDInitial\,\mathrm{m^2/s}$，对应$R^2/D(2.55)\approx\QOneMassTime\,\mathrm s$，远长于1800 s；$\sqrt{D(2.55)t}$在末时约为$\QOneMassLength\,\mathrm{mm}$，说明短时影响主要靠近表面。该尺度是扩散长度估计，不是具有锐利界面的失水前沿。
\par
\end{samepage}

此外，$D'(C)=aD(C)/C^2>0$，故失水后局部扩散系数降低。末时表面$D$约为$\QOneDLastSurface\,\mathrm{m^2/s}$，使内部向外补给的扩散能力随状态变化；这一解释直接来自题定经验式，不额外假定硬壳或相变机制。此时$C_e=\QOneCeEnd\,\mathrm{kg/kg}$，仍低于$C_s$，符合向外传质方向，也表明模型保留了表面交换阻力。

\subsection{数值验证与结果精度}

验证在相同物理时间与半径处进行，覆盖1—1800 s全部整数秒的21个径向点，并检查环境节点附近、中心、表面及短时状态。各精化解均从题定均匀初值开始，避免用粗网格后期状态掩盖启动误差的传播。零交换、平衡、常扩散圆柱解及制造解检验用于核查几何与边界实现；它们不替代真实非线性启动过程的收敛比较。

\paragraph{温度的截断与独立对照。}
模态数由512增至1024时，全部题定输出的最大差为$1.95303\times10^{-9}\,{}^{\circ}\mathrm C$。结合尾项控制及浮点余量，1024模态解的最大联合数值预算为$\QOneTEstimate\,{}^{\circ}\mathrm C$；其中解析尾界只约束截断部分，整体预算不是端到端的严格误差界。与独立2048环传热离散及其精确时间推进比较，最大差为$\QOneTIndependent\,{}^{\circ}\mathrm C$，位于1800 s、$r=1.9\,\mathrm{cm}$，该差值也包含参照的空间误差。

\paragraph{水分的空间、时间与独立离散核验。}
固定严格时间控制后，16384至32768环、32768至65536环的最大网格差分别为$4.45152\times10^{-7}$和$1.11286\times10^{-7}\,\mathrm{kg/kg}$，均在1 s表面；最大范数观测阶约为2.00003。这支持在已测网格范围内采用二阶空间估计，不表示每个局部点均严格二阶。固定32768环，Radau与BDF的最大差为$7.89862\times10^{-11}\,\mathrm{kg/kg}$；将Radau最大步长由0.5 s减至0.25 s，最大差为$1.60450\times10^{-12}\,\mathrm{kg/kg}$。这两项比较将时间误差与主要空间差异分开。

记$C_N^{\mathrm B}$为$N$环的严格BDF参照，$C_N^{\mathrm R}$为Radau结果，$C_{N,1/2}^{\mathrm R}$为最大步长减半的结果。对最终$N=32768$环全精度点值，采用以下逐点经验预算：
\begin{equation}
\widehat\varepsilon_C=
\frac{|C_N^{\mathrm B}-C_{N/2}^{\mathrm B}|}{3}
+\max\!\left\{|C_N^{\mathrm R}-C_N^{\mathrm B}|,
               |C_{N,1/2}^{\mathrm R}-C_N^{\mathrm R}|\right\}
+2\times10^{-9}\,\mathrm{kg/kg}.
\label{eq:q1-budget}
\end{equation}
第一项以观测到的近二阶行为估计细网格空间误差，包含相同点值重构的影响；第二项取时间比较的较大值，避免对同一来源重复求和；最后一项是为求解器、参考计算及浮点误差保留的经验余量，未构成严格误差界。全量最大预算为$\QOneCEstimate\,\mathrm{kg/kg}$，位于1 s表面，35个论文水分值中的最大预算为$1.29933\times10^{-8}\,\mathrm{kg/kg}$。另采用不调用主通量势和表面重构的40960区间节点离散独立核验，其自身经网格细化检查，与主解最大差为$\QOneCIndependent\,\mathrm{kg/kg}$，亦位于1 s表面。独立方法差是两种数值近似间的差异，不等于主解真实误差。

\paragraph{收支、输出与舍入检验。}
令$\bar T_i$表示环平均温度，连续模型应满足
\begin{equation}
\begin{aligned}
\frac{\mathrm d}{\mathrm dt}\sum_i\rho c_pV_i\bar T_i
 &=A_Rh_T(T_a-T_s),\\
\frac{\mathrm d}{\mathrm dt}\sum_iV_i\bar C_i
 &=-A_Rh_C(C_s-C_e).
\end{aligned}
\label{eq:q1-balance}
\end{equation}
水分累计交换量与环状态由同一时间积分器推进。全时段水分积分收支最大绝对残差为$2.5750\times10^{-19}\,\mathrm{m^3\cdot(kg/kg)}$，以初始积分量归一化后为$3.2143\times10^{-16}$。在1800 s，对独立节点剖面作Simpson积分所得的体积平均值，与主环均值积分结果相差约$2.56\times10^{-10}\,\mathrm{kg/kg}$。此处$\sum_i V_i\bar C_i$不是以kg计的实际水质量。温度原始累计能量收支残差约为$2.23\times10^{-7}\,\mathrm J$，主要由有限模态投影的截断尾项解释。这些收支检验支持实现的一致性，不能代替场值精度检验。

全部75600个四位小数结果经导出回读，与全精度权威解按约定舍入后一致；两张论文表的70项也逐项对应。经逐点分界扫描与敏感位置精化核验，现有预算下未确认项为0。同配置完整复算及65536环水分参照均支持相同舍入。其中1696 s、$r=1.1\,\mathrm{cm}$的含水率靠近分界，预算裕度较小，新增65536环参照仍支持舍入为$2.5321\,\mathrm{kg/kg}$。详细逐值证据保存在支撑材料中；这些经验数值证据不构成所有末位的严格数学保证。四位舍入本身最多改变数值约$5\times10^{-5}$，表格不保留全精度解的小误差尺度。精度结论限于题定$t\ge1\,\mathrm s$输出，亚秒启动压力测试未全部达到同一阈值。

\pagebreak[0]
\begin{samepage}
\subsection{本问结论与适用限制}

本模型给出了前30分钟内规定时刻与半径处的温度和含水率，以及完整逐秒径向结果。数值上，药材表面先升温、先失水，中心热响应滞后而含水率变化较小；这一差异与热、水分扩散时间尺度及有限表面阻力相符。方法上的特点在于利用线性热算子的半解析推进与非线性守恒通量互补，显式区分表面状态和环均值，并围绕题定输出实施分层精度核验；所用Bessel展开、有限体积与Radau均为成熟方法，不将其组合宣称为原创算法。
\par
\end{samepage}

结果的物理解释仍受两项闭合假设限制。其一，材料平衡含水率通常需要材料性质、温度与空气湿度间的关系\cite{q1-fao}；附件1未提供本药材的此类关系，将烘房水分列作为等效$C_e$的物理偏差尚未量化，不能仅凭单位一致认定二者等同。其二，蒸发会通过潜热损失影响能量平衡\cite{q1-evaporation}，本问未计入该项，现有常物性参数也不能证明其影响已被吸收。因而，本文报告的是所声明有效模型内的数值精度，未将输入有效位数、轴向近似或这些物理简化的误差一并消除。
